\chapter{Conclusions}
\label{ch:conclusion}

% PRESUPUESTO: ~3 pp

\section{Conclusions}
\label{sec:concl_conclusions}
% Volver a los TRES MODOS DE FALLO de la propuesta oficial (\ref{sec:intro_motivation})
% y tacharlos uno a uno, diciendo qué mecanismo se come cada uno.
% Cerrar el arco: el problema no estaba en el criterio, sino en la granularidad de la
% decisión y en el sustrato geométrico sobre el que se tomaba.

\section{Limitations}
\label{sec:concl_limitations}
% - El banco sintético idealiza el loop closure: corrección perfecta, instantánea, única.
% - Solo Replica. ScanNet estaba prometido en la propuesta oficial -> posicionarse.
% - El contest presupone tracking sano; fuera de ese régimen depende de los submapas.

\section{Future work}
\label{sec:concl_future}
% - Covisibilidad, si no ha entrado en \ref{sec:assoc_candidates}.
% - Señal temporal (recency/racha) en el store: aparcada, con su motivo.
%   Doc: docs/contest_temporal_signal_idea_2026-07-09.md
% - ScanNet / datos reales.
